% ============================================================
% SECTION 1 — PROJECT ORIGIN
% ============================================================

\section{Project Origin}

% --- EXECUTIVE SUMMARY BOX ---
\begin{tcolorbox}[summarybox={\iconBulb[1.5] What is RecoverAI V2?}]
\color{textdark}
RecoverAI V2 is an \textbf{autonomous revenue recovery system} designed to automatically detect, diagnose, and recover failed payment transactions. When a customer's payment fails (e.g., insufficient funds, bank timeout, card declined), RecoverAI uses a combination of \textbf{machine learning}, \textbf{LLM-powered diagnosis}, and \textbf{deterministic safety policies} to decide whether and how to retry, escalate, or contact the customer --- all without human intervention.

The system is designed with a foundational safety guarantee: \textbf{at-most-once financial execution}. A payment must never be charged twice, a refund must never be issued twice, regardless of crashes, concurrent requests, duplicate deliveries, or infrastructure failures.
\end{tcolorbox}

\vspace{0.8cm}

% --- THE PROBLEM ---
\subsection{The Problem}

\begin{tcolorbox}[infobox={\iconWarning[1.3] The Revenue Leakage Problem}]
\color{textdark}
Every online business faces a persistent problem: \textbf{payment failures}. Industry data suggests 10--15\% of online payment attempts fail. These failures have many causes:

\begin{itemize}
    \item \textbf{Insufficient funds} --- the customer's account is temporarily empty
    \item \textbf{Bank timeouts} --- the issuing bank's systems are slow or unresponsive
    \item \textbf{Card network errors} --- transient failures in Visa/Mastercard infrastructure
    \item \textbf{Fraud suspicion} --- the bank's automated fraud detection blocks the transaction
    \item \textbf{Limit exceeded} --- daily or monthly spending limits reached
\end{itemize}

Many of these failures are \textbf{recoverable}. If the system waited a few minutes and retried, or sent the customer a reminder, the payment would succeed. But most merchants simply drop these transactions, losing revenue permanently.

RecoverAI solves this by \textbf{automatically recovering failed payments} using intelligent decision-making while maintaining strict financial safety guarantees.
\end{tcolorbox}

\vspace{0.5cm}

% --- WHO IS IT FOR ---
\subsection{Target Users}

RecoverAI is designed as a \textbf{backend service} for online merchants and payment platforms. It sits between the merchant's application and the payment gateway (e.g., Razorpay, Stripe), intercepting failed transactions and attempting autonomous recovery.

\begin{multicols}{2}
\begin{tcolorbox}[featurebox={Primary Users}]
\begin{itemize}
    \item E-commerce platforms
    \item SaaS subscription services
    \item Fintech payment aggregators
    \item Any business processing online payments
\end{itemize}
\end{tcolorbox}

\columnbreak

\begin{tcolorbox}[featurebox={Integration Model}]
\begin{itemize}
    \item REST API integration
    \item Webhook-driven event processing
    \item Merchant dashboard for monitoring
    \item API key authentication
\end{itemize}
\end{tcolorbox}
\end{multicols}

\vspace{0.5cm}

% --- BUSINESS FLOW ---
\subsection{The Basic Business Flow}

In simple terms, here is what happens when a payment fails and RecoverAI processes it:

\vspace{0.3cm}

\begin{center}
\begin{tikzpicture}[
    node distance=0.8cm,
    stepbox/.style={rectangle, draw=primary!60, thick, fill=bglight, minimum width=4.5cm, minimum height=1cm, rounded corners=4pt, font=\small, text=textdark, drop shadow=black!5},
    arrow/.style={-{Stealth[scale=1.1]}, thick, draw=primary!60},
    label/.style={font=\scriptsize\itshape, text=textmuted, align=center}
]

    \node[stepbox] (cust) {Customer Attempts Payment};
    \node[stepbox, below=of cust] (fail) {Payment Fails};
    \node[stepbox, below=of fail] (ingest) {RecoverAI Ingests Transaction};
    \node[stepbox, below=of ingest] (ml) {ML Model Predicts Recovery Probability};
    \node[stepbox, below=of ml] (llm) {LLM Agent Diagnoses Failure};
    \node[stepbox, below=of llm] (policy) {Policy Engine Evaluates Safety};
    \node[stepbox, below=of policy] (guard) {ExecutionGuard Validates Invariants};
    \node[stepbox, below=of guard] (gateway) {Gateway Executes Recovery Action};
    \node[stepbox, below=of gateway] (verify) {Verification \& Final State};

    \draw[arrow] (cust) -- (fail);
    \draw[arrow] (fail) -- (ingest);
    \draw[arrow] (ingest) -- (ml);
    \draw[arrow] (ml) -- (llm);
    \draw[arrow] (llm) -- (policy);
    \draw[arrow] (policy) -- (guard);
    \draw[arrow] (guard) -- (gateway);
    \draw[arrow] (gateway) -- (verify);

    % Side labels
    \node[label, right=0.3cm of fail] {e.g., insufficient funds};
    \node[label, right=0.3cm of ml] {0.0 to 1.0 probability};
    \node[label, right=0.3cm of llm] {RETRY / WAIT / ESCALATE / STOP};
    \node[label, right=0.3cm of policy] {amount / risk / retry limits};
    \node[label, right=0.3cm of guard] {at-most-once enforcement};
    \node[label, right=0.3cm of gateway] {Razorpay / Mock};
    \node[label, right=0.3cm of verify] {SUCCEEDED / FAILED / UNKNOWN};

\end{tikzpicture}
\end{center}

\vspace{0.5cm}

% --- TECHNICAL IMPLEMENTATION BEHIND EACH STEP ---
\subsection{Technical Implementation Behind Each Step}

\begin{tcolorbox}[featurebox={Step 1: Payment Ingestion}]
\textbf{File:} \filepath{app/api/payments.py}

When a payment fails, the merchant's frontend or payment gateway sends a POST request to \code{/api/v1/payments}. The API validates the request using Pydantic schemas, checks API key authentication, enforces rate limits, and performs idempotency validation. The failed transaction is persisted to PostgreSQL, and a durable Celery task (\code{process\_orchestrator}) is enqueued for background recovery processing.
\end{tcolorbox}

\begin{tcolorbox}[featurebox={Step 2: ML Probability Prediction}]
\textbf{File:} \filepath{app/services/ml\_service.py}

A pre-trained scikit-learn model (\filepath{models/recovery\_model\_v2.pkl}) estimates the probability that this specific transaction can be recovered. The model uses hydrated features from the \code{FeatureStore} including transaction amount, currency, payment method, failure code, and historical customer behavior. The raw probability (0.0--1.0) is passed downstream without thresholding.
\end{tcolorbox}

\begin{tcolorbox}[featurebox={Step 3: LLM Diagnosis}]
\textbf{File:} \filepath{app/agents/diagnosis\_agent.py}

A configurable LLM agent (supporting OpenAI, Ollama, or a deterministic mock) analyzes the transaction context and recommends an action. The agent prompt explicitly marks all transaction data as \texttt{UNTRUSTED} to prevent prompt injection attacks. The recommended action must be one of: \code{RETRY\_PAYMENT}, \code{WAIT\_AND\_RETRY}, \code{SEND\_RECOVERY\_MESSAGE}, \code{CREATE\_ESCALATION}, or \code{STOP\_AUTOMATION}.
\end{tcolorbox}

\begin{tcolorbox}[featurebox={Step 4: Policy Engine}]
\textbf{File:} \filepath{app/policy/rules.py}

A deterministic, tiered policy engine evaluates whether the LLM's recommendation is safe to execute. It considers risk tier (based on ML probability and failure code), amount tier (SMALL/MEDIUM/LARGE), and retry count. The policy engine can \textbf{downgrade} the agent's recommendation (e.g., from RETRY to ESCALATION) but never \textbf{upgrade} it beyond the safety bounds.
\end{tcolorbox}

\begin{tcolorbox}[featurebox={Step 5: ExecutionGuard}]
\textbf{File:} \filepath{app/services/execution\_guard.py}

The final safety boundary before any financial execution. ExecutionGuard validates: (1) the action is in the allowlist, (2) the attempt exists and is in AUTHORIZED state, (3) no prior attempt for this transaction has reached a terminal or ambiguous state, (4) the transaction itself is not already terminal. Only if all invariants pass does execution proceed to the gateway.
\end{tcolorbox}

\begin{tcolorbox}[featurebox={Step 6: Gateway Execution}]
\textbf{File:} \filepath{app/services/razorpay\_mock.py}

The gateway adapter executes the actual financial operation via the payment provider's API. The current implementation uses a \code{MockGateway} that simulates Razorpay behavior, including success, failure, and timeout scenarios. Gateway-level idempotency is enforced using persistent \code{IdempotencyRecord} entries in PostgreSQL.
\end{tcolorbox}

\vspace{0.5cm}

% --- INITIAL ASSUMPTIONS ---
\subsection{Initial Assumptions}

\begin{tcolorbox}[warningbox={\iconWarning[1.3] Assumptions That Existed at the Beginning}]
\color{textdark}
When RecoverAI V2 was initially built, several assumptions were made that later proved problematic:

\begin{enumerate}
    \item \textbf{SQLite is sufficient for development and testing} --- Later discovered to hide critical concurrency issues that only manifest under PostgreSQL's true row-level locking.
    \item \textbf{In-process background tasks are reliable} --- FastAPI's \code{BackgroundTasks} drop work silently if the process crashes.
    \item \textbf{FAILED means definitively failed} --- Later discovered that broad exception handling could classify network timeouts as FAILED, making this state potentially ambiguous.
    \item \textbf{Checking \code{Transaction.recovery\_status} prevents duplicate execution} --- A crash window existed between gateway success and the transaction status update.
    \item \textbf{A single API key in source code is acceptable for development} --- Hardcoded secrets created deployment risks.
    \item \textbf{CORS wildcard (\code{*}) is acceptable} --- Created a cross-origin attack surface.
\end{enumerate}

Each of these assumptions was systematically identified and addressed across Batches 1 through 5.3.1.
\end{tcolorbox}

\vspace{0.5cm}

% --- INITIAL WEAKNESSES ---
\subsection{Initial Weaknesses}

The original RecoverAI V2 codebase had the following structural weaknesses that motivated the multi-batch hardening process:

\begin{center}
\begin{tikzpicture}[
    node distance=0.5cm,
    weak/.style={rectangle, draw=danger!50, thick, fill=danger!5, minimum width=7cm, minimum height=0.8cm, rounded corners=3pt, font=\small, text=textdark},
]
    \node[weak] (w1) {No row-level locking on financial state transitions};
    \node[weak, below=of w1] (w2) {No crash-safe execution boundary};
    \node[weak, below=of w2] (w3) {In-process background tasks (no durability)};
    \node[weak, below=of w3] (w4) {Hardcoded API keys and secrets};
    \node[weak, below=of w4] (w5) {CORS wildcard allowing any origin};
    \node[weak, below=of w5] (w6) {No webhook signature verification};
    \node[weak, below=of w6] (w7) {No request tracing or structured logging};
    \node[weak, below=of w7] (w8) {No rate limiting on financial endpoints};
    \node[weak, below=of w8] (w9) {SQLite as primary database};
    \node[weak, below=of w9] (w10) {No reconciliation for stuck refunds};
\end{tikzpicture}
\end{center}

\begin{tcolorbox}[successbox={Current Status}]
\color{textdark}
All ten of these weaknesses have been addressed across Batches 1--5.3.1. The detailed remediation history is documented in Section 4 (Batch-by-Batch History) and Section 5 (Security Evolution).
\end{tcolorbox}
